\documentclass[11pt]{article}
\usepackage[OT1]{fontenc}
\usepackage[utf8]{inputenc}
\usepackage{amsmath,amssymb,amsthm}
\usepackage[margin=1in]{geometry}
\usepackage{hyperref}
\hypersetup{hidelinks,pdftitle={Prime summand restrictions for surgery on fibered and L-space knots}}
\newtheorem{theorem}{Theorem}[section]
\newtheorem{proposition}[theorem]{Proposition}
\newtheorem{lemma}[theorem]{Lemma}
\newtheorem{corollary}[theorem]{Corollary}
\theoremstyle{remark}
\newtheorem{remark}[theorem]{Remark}
\newcommand{\Z}{\mathbb Z}
\newcommand{\Q}{\mathbb Q}
\newcommand{\F}{\mathbb F_2}
\newcommand{\HF}{\widehat{HF}}
\newcommand{\ab}{\mathrm{ab}}
\newcommand{\rk}{\operatorname{rank}}
\newcommand{\gen}{\operatorname{d}}
\title{Problem 6 -- MerLean}
\author{}
\date{}
\begin{document}
\maketitle

\begin{abstract}
We study the question whether surgery on an L-space knot has at most two prime summands. For a fibered knot of genus $g\geq1$ in $S^3$, we prove that a nonzero rational surgery whose numerator has absolute value $2g-1$ has at most two nontrivial prime summands. The proof uses the free commutator subgroup of the knot group and explicit covers of free products; it retains integral homology sphere factors that first homology alone cannot detect. A related argument shows that a nonzero integer surgery with at least three summands on a fibered hyperbolic knot must satisfy $|p|\leq2g-1$. Together with the reducing-slope theorem of Hom--Lidman--Zufelt, the first result gives the two-summand bound for hyperbolic knots in $S^3$ admitting a mod-two L-space surgery. We also give a formal obstruction to a difference-only argument and explain the role of the ambient manifold through a Seifert-fibered example. These are partial results on the original question, not an all-knot, all-coefficient conclusion.
\end{abstract}

\section{The question and the scope of the results}

The question considered here is:
\begin{quote}
Does a $3$-manifold obtained by Dehn surgery on an L-space knot have at most two prime summands?
\end{quote}
A bound on the number of reducing slopes does not itself answer this question: a single filling could have several prime factors. Moreover, the cyclic first homology of surgery on a knot in $S^3$ does not see integral homology sphere factors. Our purpose is to give group-theoretic restrictions that address these two issues and to specify the cases to which they apply.

All manifolds are smooth, connected and oriented. A closed $3$-manifold is \emph{prime} if every connected-sum decomposition has an $S^3$ factor. We count the non-$S^3$ factors in its prime decomposition, with multiplicity; in particular, $S^1\times S^2$ counts as one factor and $S^3$ counts as zero. A sphere is \emph{essential} if it does not bound a ball, and a manifold is \emph{reducible} if it contains such a sphere. We use the prime decomposition theorem and the Poincar\'e theorem; consequently every counted prime factor has nontrivial fundamental group. For background on prime decomposition see \cite[Section~1.1]{Hatcher}.

Unless explicitly stated otherwise, the knot lies in $S^3$. Write
\[
 E(K)=S^3\setminus\operatorname{int}\nu K,
 \qquad S^3_{p/q}(K)=E(K)(p\mu+q\lambda),
\]
where $\nu K$ is a tubular neighborhood, $\mu$ is a meridian and $\lambda$ is the preferred longitude. A finite slope has coprime integers $p,q$ with $q>0$; the meridian is $\infty=1/0$. Dehn filling identifies the indicated slope with the meridian of the attached solid torus. The Seifert genus is denoted by $g(K)$. A knot is \emph{fibered} if its exterior is a bundle over $S^1$ with fiber a connected Seifert surface. It is \emph{hyperbolic} if the interior of its exterior has a complete finite-volume hyperbolic metric.

For the Floer consequence we fix coefficients: an $\F$-\emph{L-space} is a rational homology sphere $Y$ satisfying
\[
 \dim_{\F}\HF(Y;\F)=|H_1(Y;\Z)|.
\]
Here $\HF$ denotes hat Heegaard Floer homology. We say that $K$ admits an $\F$-L-space surgery if some finite nonzero rational filling is an $\F$-L-space. Excluding the meridian is necessary, since meridional surgery on every knot in $S^3$ returns $S^3$. The surgery being tested in the question is independent of this witnessing surgery; it is not assumed to be an L-space.

\begin{theorem}\label{thm:main}
The following statements hold.
\begin{enumerate}
\item If $K\subset S^3$ is fibered of genus $g\geq1$ and $p/q$ is reduced with $|p|=2g-1$, then $S^3_{p/q}(K)$ has at most two counted prime factors.
\item If $K\subset S^3$ is fibered and hyperbolic of genus $g$, and a nonzero integer surgery $S^3_p(K)$ has at least three counted prime factors, then $|p|\leq2g-1$.
\item If $K\subset S^3$ is hyperbolic and admits an $\F$-L-space surgery, then every filling $S^3_r(K)$, $r\in\Q\cup\{\infty\}$, has at most two counted prime factors.
\end{enumerate}
\end{theorem}

The first two statements do not use Heegaard Floer homology. Their central input is that the commutator subgroup of a fibered knot group is free of rank $2g$. We compare its quotients with the commutator groups forced by a connected sum. The third statement uses Ni's fiberedness theorem \cite[Corollary~1.3]{Ni}, with its correction \cite{NiErr}, and Hom--Lidman--Zufelt's reducing-slope theorem \cite[Theorem~1.3]{HLZ}. Their coefficient and sign conventions are checked in Section~\ref{sec:floer}.

There is relevant earlier work on both parts of this argument. Bodish--DeYeso \cite[Theorem~1.2]{BD} record a genus bound for a fibered knot surgery of the form $L(r,a)\#Y$, with $Y$ a homology sphere, and attribute it to Gonz\'alez-Acu\~na--Short. We give a direct free-product proof of the estimate needed here, allowing an arbitrary complementary summand; no lens-space identification is required. The group-theoretic input is the Grushko--Wagner theorem, in the free-source form stated by Stallings \cite[Section~1.B.5, p.~6]{Stallings}. Bodish--DeYeso also use differences of correction terms to obstruct reducible surgery. Section~\ref{sec:formal} gives a precise limitation of one purely algebraic version of that method. We make no priority claim for the group estimates.

Theorem~\ref{thm:main} does not identify every possible meaning of ``L-space knot,'' nor does it supply the full satellite-knot reduction needed for an all-knot answer. Section~\ref{sec:limits} distinguishes the remaining issues from the results proved here. The broader ambient example there concerns a different statement and is not a counterexample for knots in $S^3$.

\section{Fibered knot groups and invisible homology factors}\label{sec:groups}

For a group $A$, write $A'=[A,A]$ and $A_{\ab}=A/A'$. Let $F_k$ be the free group of rank $k$, let $C_m=\Z/m\Z$ for $m\geq1$, and let $\gen(A)$ be the least number of generators of a finitely generated group $A$. An epimorphism cannot increase this number. We will not use any monotonicity assertion for arbitrary subgroups.

The Grushko--Wagner theorem says that an epimorphism from a free group onto a free product admits a decomposition of the source into free factors mapping onto the respective target factors \cite[Section~1.B.5]{Stallings}. Applied to a minimal finite free source, it gives
\begin{equation}\label{eq:grushko}
 \gen(A_1*\cdots*A_n)=\sum_{i=1}^n\gen(A_i)
\end{equation}
for finitely generated groups. The opposite inequality follows by taking the union of generating sets. This is the external group theorem used below.

\begin{lemma}\label{lem:fiber}
Let $K\subset S^3$ be fibered of genus $g$. For every filling $M$ with fundamental group $G$, there is an epimorphism
\[
 F_{2g}\longrightarrow G'.
\]
For reduced $p/q$ with $p\ne0$, one also has $H_1(S^3_{p/q}(K);\Z)\cong C_{|p|}$.
\end{lemma}
\begin{proof}
The exterior has first homology $\Z$, generated by $\mu$, while $\lambda$ is zero in this group. Mayer--Vietoris for the filling imposes the relation $p\mu=0$, giving the stated homology calculation. For $p=0$ the same calculation gives $\Z$.

Put $\Gamma=\pi_1(E(K))$. The fibration induces an epimorphism $\Gamma\to\Z$ whose kernel is the fundamental group of the once-punctured genus-$g$ fiber, namely $F_{2g}$. Since $\Gamma_{\ab}=\Z$, this epimorphism is abelianization up to sign. Thus $\Gamma'=F_{2g}$. The filling map $\Gamma\to G$ is onto. Every commutator in $G$ is the image of a commutator in $\Gamma$, by lifting its two entries; the same is true of products of commutators. Hence $\Gamma'\to G'$ is onto.
\end{proof}

\begin{lemma}\label{lem:perfect-cover}
Suppose $A_{\ab}=C_m$ and $B$ is perfect, meaning $B'=B$. Then
\[
 (A*B)'\cong A'*\underbrace{B*\cdots*B}_{m\text{ copies}}.
\]
\end{lemma}
\begin{proof}
Take connected presentation complexes $X_A,X_B$, each with one vertex, and form their wedge. Its fundamental group is $A*B$. Consider the connected cover corresponding to the kernel of the abelianization map $A*B\to C_m$. The preimage of $X_A$ is one connected $m$-sheeted cover with fundamental group $A'$, because the restriction $A\to C_m$ is onto. The preimage of $X_B$ consists of $m$ copies of $X_B$, each attached at a single lifted vertex, since $B$ maps trivially to $C_m$.

Choose paths in the connected $A$-cover from a base point to these vertices. Van Kampen now gives one free factor $A'$ and one free factor $B$ for each lifted vertex. Attaching a space at one point introduces no additional graph cycle, so there is no extra free factor. The kernel defining the cover is exactly $(A*B)'$.
\end{proof}

\begin{proposition}\label{prop:homology-sphere}
Let $K\subset S^3$ be fibered of genus $g$, and let $p/q$ be reduced with $p\ne0$. If
\[
 S^3_{p/q}(K)=X\#Y
\]
where $Y\ne S^3$ is an integral homology sphere, then $|p|\leq g$.
\end{proposition}
\begin{proof}
Write $m=|p|$, $A=\pi_1(X)$ and $B=\pi_1(Y)$. Connected-sum homology is a direct sum, so $A_{\ab}=C_m$ and $B_{\ab}=0$. The group $B$ is finitely generated by compactness and nontrivial by the Poincar\'e theorem. A nontrivial perfect group cannot be cyclic, and therefore $\gen(B)\geq2$.

Van Kampen, Lemma~\ref{lem:perfect-cover}, and then the quotient killing $A'$ give epimorphisms
\[
 F_{2g}\longrightarrow (A*B)'\longrightarrow B^{*m}.
\]
Equation~\eqref{eq:grushko} implies
\[
 2g\geq\gen(B^{*m})=m\gen(B)\geq2m.
\]
No finite-generation assumption on $A'$ is needed.
\end{proof}

This estimate explains why an integral homology sphere factor cannot simply be discarded. Although its abelianization vanishes, the abelianization cover of the entire filling contains $|p|$ copies of its fundamental group.

\section{Cyclic quotients and the odd endpoint}\label{sec:endpoint}

Suppose a nonzero filling has a counted prime decomposition
\begin{equation}\label{eq:sum}
 M=M_1\#\cdots\#M_n,
 \qquad A_i=\pi_1(M_i),\qquad m=|p|.
\end{equation}
Since $\bigoplus_i(A_i)_{\ab}=C_m$, each factor is cyclic, say $(A_i)_{\ab}=C_{a_i}$, and
\begin{equation}\label{eq:orders}
 \prod_i a_i=m,\qquad \gcd(a_i,a_j)=1\quad(i\ne j).
\end{equation}
Indeed, each direct summand embeds in $C_m$, and two subgroups of a cyclic group of orders $a_i,a_j$ intersect in order $\gcd(a_i,a_j)$. Distinct direct summands intersect trivially. If $a_i=1$, Poincar\'e duality gives $H_2(M_i;\Z)=H^1(M_i;\Z)=0$, so $M_i$ is an integral homology sphere. Cyclic homology alone does not identify $A_i$ as a cyclic group or $M_i$ as a lens space.

When every $a_i>1$, abelianizing the individual factors gives an epimorphism
\[
 \pi_1(M)=*_{i=1}^n A_i\longrightarrow H:=*_{i=1}^n C_{a_i}.
\]
We compute $H'$ explicitly to obtain a lower bound on the number of generators needed for $\pi_1(M)'$.

\begin{lemma}\label{lem:cyclic-cover}
For integers $a_i\geq2$ and $m=\prod_i a_i$, the commutator subgroup of $H=*_{i=1}^n C_{a_i}$ is free of rank
\begin{equation}\label{eq:R}
 R=1+(n-1)m-\sum_{i=1}^n\frac{m}{a_i}.
\end{equation}
Pairwise coprimality is not required for this formula.
\end{lemma}
\begin{proof}
Use the presentation $\langle x_1,\ldots,x_n\mid x_i^{a_i}=1\rangle$. The cover corresponding to $H'$ has vertex set $Q=\prod_i C_{a_i}$ and $nm$ directed edge cells
\[
 (v,i):v\longrightarrow v+e_i,
\]
where $e_i$ is the standard generator of the $i$th cyclic factor. For each fixed $i$, these edges form $m/a_i$ disjoint cycles of length $a_i$. If $a_i=2$, the two directed edges in such a cycle remain distinct cells.

Every lift of the $i$th relator traverses one of these cycles once. The different starting vertices on that cycle give conjugate attaching loops. Retain just one disk on each cycle. Removing the other disks leaves the fundamental group unchanged, because their attaching loops are already trivial after the retained disk is attached. This is an assertion about fundamental groups, not a deformation retraction of the original complex.

In the resulting complex collapse each retained disk across one edge in its boundary. Each chosen edge occurs once in that boundary and belongs to no other retained disk: the cycle edge sets are disjoint for a fixed $i$, and different $i$ use differently labeled edges. The resulting graph is connected, since a removed edge can be replaced by the rest of its cycle. It has $m$ vertices and
\[
 nm-\sum_i m/a_i
\]
edges. Its fundamental group is free of rank the number of edges minus the number of vertices plus one, proving \eqref{eq:R}.
\end{proof}

\begin{lemma}\label{lem:odd}
If $n\geq3$ and the $a_i>1$ are odd and pairwise coprime, then $R>m+1$.
\end{lemma}
\begin{proof}
It suffices to show $\sum_i1/a_i<n-2$. For $n=3$, each $a_i\geq3$ gives a sum at most one; equality would require all three to equal $3$, contradicting coprimality. For $n>3$, the sum is at most $n/3<n-2$.
\end{proof}

\begin{proof}[Proof of Theorem~\ref{thm:main}(1)]
Suppose $n\geq3$ in \eqref{eq:sum}, and put $m=|p|=2g-1$. By Lemma~\ref{lem:fiber}, $F_{m+1}$ surjects onto $G'=\pi_1(M)'$.

If $m=1$, every $A_i$ is nontrivial and perfect. Then $G'=G=*A_i$, and Grushko gives
\[
 \gen(G')=\sum_i\gen(A_i)\geq2n\geq6>2=m+1,
\]
a contradiction. If $m>1$ and some $a_i=1$, the proof of Proposition~\ref{prop:homology-sphere} gives $\gen(G')\geq2m>m+1$, again a contradiction.

Otherwise all $a_i>1$. They are odd and pairwise coprime by \eqref{eq:orders}. The factorwise abelianization induces an epimorphism $G'\to H'$, and Lemmas~\ref{lem:cyclic-cover} and \ref{lem:odd} give $H'=F_R$ with $R>m+1$. Thus there would be an epimorphism $F_{m+1}\to F_R$. Abelianizing and tensoring with $\Q$ would give an epimorphism $\Q^{m+1}\to\Q^R$, which is impossible.
\end{proof}

For illustration, orders $(3,5,7)$ give $m=105$ and $R=140$, whereas $(3,5)$ give $m=15$ and $R=8$. These are calculations of abstract free-product ranks, not surgery examples. Oddness matters: $(2,3,5)$ give $m=R=30$. The next section deals with precisely this exceptional configuration.

\section{A genus bound when three summands are assumed}\label{sec:genus}

The odd endpoint argument does not give a uniform inequality for even $m$. Nevertheless, a three-summand hypothesis yields enough control except at one equality case. Hyperbolicity excludes that equality case by detecting the kernel of the actual filling map.

\begin{lemma}\label{lem:recip}
Let $n\geq3$, with pairwise coprime $a_i>1$. In the notation of Lemma~\ref{lem:cyclic-cover}, $R>m+1$ unless $n=3$ and the unordered orders are $\{2,3,5\}$. In that case $m=R=30$.
\end{lemma}
\begin{proof}
For $n\geq4$, at most one order is $2$, and all others are at least $3$. Hence
\[
 \sum_i\frac1{a_i}\leq\frac12+\frac{n-1}{3}<n-2,
\]
the last difference being $(4n-13)/6>0$.

For $n=3$, order the integers as $a\leq b\leq c$. If $a\geq3$, their reciprocal sum is at most one, with equality only when all equal $3$, which is excluded. If $a=2$ and $b\geq5$, then $b$ is odd and $c\geq7$, so the sum is at most $1/2+1/5+1/7<1$. The remaining case has $a=2,b=3$. Coprimality gives either $c=5$ or $c\geq7$; in the latter case the sum is at most $41/42<1$. For $(2,3,5)$, formula~\eqref{eq:R} gives $1+60-15-10-6=30$.
\end{proof}

We include the group facts needed to use this arithmetic equality, rather than inferring injectivity from equality of generator numbers for arbitrary groups.

\begin{lemma}\label{lem:hopf}
An epimorphism between free groups of the same finite rank is an isomorphism.
\end{lemma}
\begin{proof}
First, $F_k$ is residually finite: every nonidentity element survives under a homomorphism to a finite group. For a nonempty reduced word $w$ of length $L$, form a path with vertices $0,\ldots,L$ and successive edge labels the letters of $w$. A letter $x$ at position $j$ prescribes $x:j-1\mapsto j$, and $x^{-1}$ prescribes $x:j\mapsto j-1$. For each free generator these prescriptions form a partial injection. Conflicting sources or targets could only occur at a shared vertex of adjacent edges, and would require adjacent inverse letters, contrary to reduction. Complete each partial injection to a permutation of this finite vertex set. With the right-action convention, $w$ sends $0$ to $L$, so it has nontrivial image.

A finitely generated residually finite group $J$ is Hopfian, meaning that every onto endomorphism is injective. To see this, let $f:J\to J$ be onto. For any finite group $Q$, the set $\operatorname{Hom}(J,Q)$ is finite. Precomposition by $f$ is injective on this set and hence onto. If $1\ne w\in\ker f$, choose $q:J\to Q$ with $q(w)\ne1$. Surjectivity of precomposition writes $q=h\circ f$, a contradiction. Apply this to $F_k$, after choosing an isomorphism between the source and target free groups.
\end{proof}

\begin{lemma}\label{lem:center}
If an epimorphism $\phi:\Gamma\to G$ is injective on $\Gamma'$, then $\ker\phi$ is central in $\Gamma$. The group of a fibered hyperbolic knot in $S^3$ has trivial center.
\end{lemma}
\begin{proof}
For $z\in\ker\phi$ and $\gamma\in\Gamma$, the commutator $[\gamma,z]$ belongs to both $\ker\phi$ and $\Gamma'$. Their intersection is trivial, proving the first assertion.

A fibered hyperbolic knot has genus $g\geq1$. Its group $\Gamma$ is a faithful torsion-free subgroup of $\operatorname{PSL}_2(\mathbb C)$ and is nonabelian, since $\Gamma'=F_{2g}$. Suppose $1\ne z$ is central, and lift it to a nonscalar matrix $Z\in\operatorname{SL}_2(\mathbb C)$. Its trace is nonzero: if it were zero, Cayley--Hamilton would give $Z^2=-I$, making $z$ an element of order two. Every lift $A$ of an element of $\Gamma$ satisfies $AZA^{-1}=\pm Z$; equality of traces excludes the minus sign. The centralizer of a nonscalar complex $2\times2$ matrix is abelian. For distinct eigenvalues this follows by diagonalization, and for a single nonscalar Jordan block its centralizer consists of matrices $aI+bN$ in a Jordan basis. Thus all lifts of $\Gamma$ commute, a contradiction.
\end{proof}

\begin{proof}[Proof of Theorem~\ref{thm:main}(2)]
A hyperbolic fibered knot has $g\geq1$. Put $m=|p|$ and suppose, to the contrary, that $m\geq2g$ and the filling has $n\geq3$ counted prime factors. Use the actual groups and homology orders in \eqref{eq:sum}--\eqref{eq:orders}.

If some $a_i=1$, Proposition~\ref{prop:homology-sphere} gives $2g\geq2m$, contradicting $m\geq2g$ and $g\geq1$. Otherwise factorwise abelianization and Lemmas~\ref{lem:fiber} and \ref{lem:cyclic-cover} yield epimorphisms
\[
 \Gamma'=F_{2g}\longrightarrow G'\longrightarrow H'=F_R.
\]
Consequently $2g\geq R$. Lemma~\ref{lem:recip} contradicts this unless the orders are $(2,3,5)$. In that case $m=R=30$, and the inequalities $2g\geq R$ and $m\geq2g$ force $2g=30$.

The composite $F_{30}\to G'\to F_{30}$ is now injective by Lemma~\ref{lem:hopf}; therefore its first map is injective. By Lemma~\ref{lem:center}, the kernel of the actual filling epimorphism $\Gamma\to G$ is central. But that kernel contains $\mu^p\lambda$, whose image in $\Gamma_{\ab}=\Z$ is the nonzero integer $p$. Thus the kernel contains a nonidentity element. This contradicts the trivial center and excludes the last case. Hence $|p|\leq2g-1$.
\end{proof}

\section{Application to the hyperbolic L-space case}\label{sec:floer}

Two published Floer-theoretic results supply the exact topological hypotheses for Theorem~\ref{thm:main}(1). Ni's Corollary~1.3 \cite{Ni}, with the proof correction in \cite{NiErr}, states that a knot in $S^3$ with an L-space surgery is fibered. Hom--Lidman--Zufelt work throughout over $\F$ and define an L-space knot using a positive surgery witness \cite[Section~1, p.~2]{HLZ}. Their Theorem~1.3 states that for a hyperbolic L-space knot every reducing slope equals $2g(K)-1$. The tested filling need not itself be an L-space. We use this theorem for exactly that coefficient convention.

For clarity, the change from mod-two minimal rank to the integral-rank convention stated immediately before Ni's corollary is one-way. Let $C$ be the finite free integral hat Floer complex of the witnessing rational homology sphere $Y$. Coherent integral orientations give the integral complex whose reduction is the mod-two complex. Universal coefficients and the Euler-characteristic lower bound give
\begin{equation}\label{eq:UCT}
 |H_1(Y;\Z)|\leq\rk_{\Z}H(C)
 \leq\dim_{\F}H(C\otimes\F).
\end{equation}
For an $\F$-L-space witness equality holds throughout, meeting the rank condition in Ni's corollary. This calculation does not assert that $H(C)$ is torsion-free. Alternatively, the fiberedness of knots with the mod-two convention is recorded explicitly in \cite[Section~1]{HLZ}.

\begin{proof}[Proof of Theorem~\ref{thm:main}(3)]
Reflection identifies $S^3_r(K)$ with the orientation reverse of $S^3_{-r}(\overline K)$, where $\overline K$ is the mirror. It preserves prime count, and Floer orientation duality preserves total rank over $\F$. Negation permutes all tested slopes and fixes $0$ and $\infty$. We may therefore normalize the witnessing slope to be positive. This does not assume the tested slope is positive.

Ni's theorem gives fiberedness. Meridional surgery is $S^3$. Suppose some other filling has at least three counted prime factors. A connected-sum sphere separating one prime factor from the rest is essential: if either punctured side were a ball, its cap would be $S^3$, contrary to uniqueness of the nontrivial prime decomposition. The filling is therefore reducible. The theorem of Hom--Lidman--Zufelt applies to the hyperbolic knot with its positive $\F$-L-space witness and makes this slope the integer $2g-1$. Theorem~\ref{thm:main}(1) rules out three factors at that slope.
\end{proof}

Some elementary cases require no Floer theory. The meridional filling of any knot is $S^3$. Fillings of the unknot are $S^3$, $S^1\times S^2$, or lens spaces, and so have at most one counted prime factor. Here a lens space means a manifold obtained by gluing two solid tori; when its first homology has finite order greater than one it is a nontrivial prime manifold.

For a torus knot with relatively prime winding numbers $u,v$, van Kampen gives the exterior group $\langle a,b\mid a^u=b^v\rangle$. Every filling group is a quotient and hence is at most two-generated. Three counted prime factors would have three nontrivial groups and, by Grushko, would require at least three generators. Thus all torus-knot fillings also satisfy the bound. This argument uses only the exterior presentation, not a classification of all its fillings.

\section{A limitation of a formal difference obstruction}\label{sec:formal}

The following calculation concerns rational arrays, not realized correction terms. Its relevance is that additive connected-sum formulas motivate looking for a decomposition of a function on a cyclic group into functions of two residue classes; compare the correction-term methods in \cite[Section~2.1]{BD}. Such a decomposition forces vanishing mixed differences. The converse information is too weak for the data below.

For a Laurent polynomial $Q(t)=\sum_jq_jt^j$ and $s\geq0$, put
\[
 w_s(Q)=\sum_{j>s}(j-s)q_j.
\]
Define the formal polynomials
\begin{align*}
 A(t)&=t^4-t^3+t-1+t^{-1}-t^{-3}+t^{-4},\\
 P(t)&=A(t)+t^8-t^7+t^{-8}-t^{-7}.
\end{align*}
The second polynomial is symmetric, has $P(1)=1$, and has alternating nonzero coefficients. Direct coefficient sums give, for $s=0,\ldots,8$,
\begin{align*}
 (w_s(A))&=(2,1,1,1,0,0,0,0,0),\\
 (w_s(P))&=(3,2,2,2,1,1,1,1,0).
\end{align*}
Indeed, the two added positive powers contribute $(8-s)-(7-s)=1$ for $0\leq s\leq6$, contribute $1$ at $s=7$, and contribute zero thereafter. The negative powers do not enter these sums.

The descending exponents of $P$ are
\[
 (8,7,4,3,1,0,-1,-3,-4,-7,-8).
\]
Their successive gaps give a formal staircase by alternating rightward and downward moves from $(0,8)$. Its vertices are
\[
 \begin{gathered}
 (0,8),(1,8),(1,5),(2,5),(2,3),(3,3),\\
 (3,2),(5,2),(5,1),(8,1),(8,0).
 \end{gathered}
\]
Here ``formal staircase'' means only this polygonal path. Number its vertices from zero. The minimum of $\max(x,y-s)$ over its even-numbered vertices equals $w_s(P)$ for $s=0,\ldots,8$: the six quantities being minimized are
\[
 \begin{gathered}
 \max(0,8-s),\quad \max(1,5-s),\quad \max(2,3-s),\\
 \max(3,2-s),\quad \max(5,1-s),\quad 8.
 \end{gathered}
\]
For every $s\geq8$ the first vertex gives the minimum zero. No knot Floer complex is asserted to realize this path.

For $0\leq i<15$, define
\[
 D_Q(i)=\frac{(2i-15)^2-15}{60}-2w_{\min(i,15-i)}(Q).
\]
All indices in the second term are at most $7$, so $D_P(i)=D_A(i)-2$. With arrays indexed from zero, put
\[
 F=(1/6,1/6,-1/2),\qquad
 G=(2/5,0,2/5,-2/5,-2/5).
\]
Then every one of the fifteen values satisfies
\begin{equation}\label{eq:formal-factor}
 D_P(i)=F_{(i+2)\bmod3}+G_{(i+1)\bmod5}-2.
\end{equation}
For an explicit check, the two sides of \eqref{eq:formal-factor}, multiplied by $30$ and listed in order, both give
\[
 \begin{gathered}
 (-75,-43,-67,-87,-43,-55,-63,-67,\\
 -67,-63,-55,-43,-87,-67,-43).
 \end{gathered}
\]
Extend this function periodically to $\Z$. Its mixed difference with shifts $3$ and $5$ vanishes, because in \eqref{eq:formal-factor} the $F$ terms cancel under the shift by $3$ and the $G$ terms cancel under the shift by $5$.

Thus symmetry, alternating coefficients, the displayed staircase test, and this residue-class factorization do not contradict each other. This is not a counterexample to the surgery question: it supplies neither an L-space knot realizing $P$, nor a surgery decomposition, nor an identification of these arrays with actual labeled correction terms. It only rules out an obstruction based solely on the stated formal conditions.

\section{Ambient and coefficient limitations}\label{sec:limits}

The ambient $S^3$ assumption enters through the abelianization of the exterior and the identification of the fiber subgroup with the commutator subgroup. Allowing knots in other L-space integral homology spheres changes the question. The following construction makes that distinction concrete.

\begin{proposition}\label{prop:ambient}
There is a prime integral homology sphere $Y\ne S^3$ which is an L-space, and a nullhomologous knot $K\subset Y$, such that $31$-surgery on $K$ is an L-space and $30$-surgery is a connected sum of exactly three nontrivial lens spaces. The exterior of $K$ is not homeomorphic to the exterior of any knot in $S^3$.
\end{proposition}
\begin{proof}
Let $P_0$ be a disk with three open disks removed. Orient its inner boundary curves $x_1,x_2,x_3$ and outer boundary curve $d$ so that $d+x_1+x_2+x_3=0$ in homology. Let $h$ be the $S^1$ factor of $P_0\times S^1$. Fill the three inner boundary tori along
\[
 2x_1+h,\qquad3x_2+h,\qquad5x_3+h,
\]
and call the remaining torus-boundary manifold $W$. This is the usual construction of a Seifert fibration over a disk with exceptional multiplicities $2,3,5$: a filling meridian meeting a regular fiber $a$ times creates multiplicity $a$; see \cite[Section~2.1, pp.~24--26]{Hatcher}.

The homology presentation has generators $x_1,x_2,x_3,h$ and relations $2x_1+h=3x_2+h=5x_3+h=0$. Its relation matrix has maximal minors with absolute values $30,15,10,6$, whose greatest common divisor is one. Its rank is three, so $H_1(W;\Z)=\Z$. An identification with $\Z$ is
\[
 (x_1,x_2,x_3,h)\longmapsto(-15,-10,-6,30),\qquad d\longmapsto31.
\]
Set
\[
 \mu=d-h,\qquad \lambda=h-30\mu=-30d+31h.
\]
These form a basis of the boundary homology, since their determinant in the $(d,h)$ basis is one. Moreover, $\mu$ generates $H_1(W)$ and $\lambda$ vanishes there. Thus $Y=W(\mu)$ is an integral homology sphere, and the core $K$ of this filling is a nullhomologous knot with exterior $W$, meridian $\mu$ and preferred longitude $\lambda$. On this same knot,
\[
 h=30\mu+\lambda,\qquad d=31\mu+\lambda.
\]

Both $\mu$ and $d$ meet $h$ once. Their fillings have base orbifold $S^2(2,3,5)$, with Euler numbers, up to an overall orientation sign,
\[
 e(Y)=-1+\frac12+\frac13+\frac15=\frac1{30},
 \qquad e(W(d))=\frac12+\frac13+\frac15=\frac{31}{30}.
\]
The base orbifold has Euler characteristic $-1+1/2+1/3+1/5=1/30>0$. A Seifert manifold with this positive orbifold Euler characteristic and nonzero Euler number has elliptic geometry; this criterion and the resulting L-space assertion are given in \cite[Proposition~2.3 and its proof]{OSLens}. Hence $Y$ and $W(d)$ are integral L-spaces in the free-abelian convention of \cite[Definition~1.1]{OSLens}, and therefore are also $\F$-L-spaces by universal coefficients. An elliptic $3$-manifold is irreducible, so $Y$ is prime.

To see the three factors of $W(h)$ geometrically, take disjoint arcs in the base disk which separate its three exceptional points. The inverse images are vertical annuli with boundary slope $h$. In $W(h)$ their boundary circles cap by disjoint meridian disks to separating spheres. Cutting along an annulus divides the base into disks and cuts the attached filling solid torus into balls. After capping a resulting sphere by a ball, the portion of the filling torus together with that cap is a solid torus completing the regular-fiber filling on that side. Iterating gives three pieces, each a solid torus with one exceptional fiber, filled along its regular fiber. In such a piece the regular fiber represents $a$ times the core, with $a=2,3,5$, respectively. Gluing in the filling solid torus therefore gives a lens space with first homology of order $a$. Thus
\[
 W(h)=L_2\#L_3\#L_5,\qquad |H_1(L_j;\Z)|=j.
\]
These are three nontrivial prime factors. This splitting is not inferred from first homology alone.

It remains to check that the construction cannot be reinterpreted as an $S^3$ knot exterior. Van Kampen gives
\[
 \begin{gathered}
 \pi_1(W)=\langle x_1,x_2,x_3,h\mid [x_i,h]=1,\
 x_1^2h=x_2^3h=x_3^5h=1\rangle,\\
 d=(x_1x_2x_3)^{-1}.
 \end{gathered}
\]
In the alternating group $A_5$, with permutations composed right to left, take
\[
 a=(12)(34),\qquad b=(135),\qquad c=(ab)^{-1}.
\]
Then $ab=(1\ 4\ 3\ 5\ 2)$, so $a^2=b^3=c^5=abc=1$. Sending $(x_1,x_2,x_3,h)$ to $(a,b,c,1)$ defines an epimorphism onto $A_5$. Indeed its image has order divisible by $2,3,5$ and divides $60$. An image of order $30$ would be an index-two subgroup of $A_5$, giving a nontrivial homomorphism $A_5\to C_2$, which is impossible because $A_5$ is generated by $3$-cycles. For completeness, pairs of transpositions sharing a letter give $3$-cycles, while disjoint pairs satisfy $(ab)(cd)=(a\ c\ b)(a\ c\ d)$; thus every even permutation is a product of $3$-cycles.

Both peripheral generators $d,h$ map to the identity, so this epimorphism factors through every Dehn filling of $W$. No filling can be $S^3$, and in particular $Y\ne S^3$. If $W$ were homeomorphic, even without a chosen boundary basis, to a knot exterior in $S^3$, the inverse image of its meridional slope would give an $S^3$ filling. This contradiction proves the last assertion.
\end{proof}

The connected sum $L_2\#L_3\#L_5$ is itself an integral L-space: lens spaces are integral L-spaces and that class is closed under connected sum \cite[paragraph following Definition~1.1]{OSLens}. Thus even requiring the tested filling to be an L-space does not repair the broader ambient statement. Proposition~\ref{prop:ambient} does not change the domain of Theorem~\ref{thm:main}.

There is also a coefficient distinction. For a rational homology sphere with an actual finite free integral hat complex $C$, flatness gives
\[
 \dim_{\Q}H(C\otimes\Q)=\rk_{\Z}H(C).
\]
For a field $k$ of positive characteristic, universal coefficients followed by scalar extension gives $\dim_kH(C\otimes k)\geq\rk_{\Z}H(C)$. Together with the Euler-characteristic lower bound, minimal rank over any field implies minimal rank over $\Q$ at the same slope. An integral-free witness likewise gives a mod-two witness. These are one-way implications, not equivalences.

For example, the complex $C_1=\Z$, $C_0=\Z^2$ with differential $x\mapsto(2x,0)$ has integral homology $H_0=\Z\oplus\Z/2$ and $H_1=0$, but total mod-two homology dimension three. Replacing $2$ by $3$ gives the same integral rank and total mod-two dimension one. These are abstract complexes, not asserted Floer complexes of manifolds. They explain why rational minimal rank alone does not justify applying a theorem stated over $\F$.

For a rational-rank witness, extending the hyperbolic conclusion would require a theorem about the actual surgery Floer groups over $\Q$. One sufficient missing statement is the following endpoint assertion: if a nontrivial hyperbolic knot with a rational-rank L-space witness has a nonzero integer surgery with a nontrivial lens-space summand and $|p|\leq2g-1$, then $|p|=2g-1$. Theorem~\ref{thm:main}(2) supplies this magnitude bound when three factors are assumed; it supplies neither the lens summand nor the coefficient-specific endpoint assertion. The latter is not proved in this paper. Similarly, a complete all-knot conclusion needs the reducible satellite-filling reduction and its peripheral identifications. The local estimates and formal examples above do not replace those steps. Accordingly, the original question remains unanswered here beyond the cases explicitly proved.

\begin{thebibliography}{9}
\bibitem{BD}
H.~Bodish and R.~DeYeso~III,
\emph{Obstructing reducible surgeries: slice genus and thickness bounds},
arXiv:2209.01672v1 (2022).
\href{https://arxiv.org/abs/2209.01672v1}{arXiv:2209.01672v1}.

\bibitem{Hatcher}
A.~Hatcher, \emph{Notes on Basic 3-Manifold Topology},
author's notes, revisions 1999--2000.
\href{https://pi.math.cornell.edu/~hatcher/3M/3M.pdf}{Author's online text}.

\bibitem{HLZ}
J.~Hom, T.~Lidman and N.~Zufelt,
\emph{Reducible surgeries and Heegaard Floer homology},
Math. Res. Lett. \textbf{22} (2015), no.~3, 763--788.
\href{https://arxiv.org/abs/1307.5317v2}{arXiv:1307.5317v2}.

\bibitem{Ni}
Y.~Ni, \emph{Knot Floer homology detects fibred knots},
Invent. Math. \textbf{170} (2007), no.~3, 577--608.
\href{https://doi.org/10.1007/s00222-007-0075-9}{doi:10.1007/s00222-007-0075-9}.
\href{https://arxiv.org/abs/math/0607156v4}{arXiv:math/0607156v4}.

\bibitem{NiErr}
Y.~Ni, \emph{Erratum: Knot Floer homology detects fibred knots},
Invent. Math. \textbf{177} (2009), 235--238.
\href{https://doi.org/10.1007/s00222-009-0174-x}{doi:10.1007/s00222-009-0174-x}.

\bibitem{OSLens}
P.~Ozsv\'ath and Z.~Szab\'o,
\emph{On knot Floer homology and lens space surgeries},
Topology \textbf{44} (2005), 1281--1300.
\href{https://arxiv.org/abs/math/0303017v2}{arXiv:math/0303017v2}.

\bibitem{Stallings}
J.~Stallings, \emph{Group Theory and Three-Dimensional Manifolds},
Yale Mathematical Monographs, vol.~4, Yale University Press,
New Haven and London, 1971.
\href{https://webhomes.maths.ed.ac.uk/~v1ranick/papers/stallings1.pdf}{Online monograph scan}.
\end{thebibliography}
\end{document}
